% #############################################################################
% This is Chapter 3
% !TEX root = ../main.tex
% #############################################################################
% Change the Name of the Chapter i the following line
\fancychapter{A CLASS Compiler}
% The following line allows to ref this chapter
\label{chap:architecture}

The solution developed on this thesis to the problem of compiling \ac{CLASS} programs to native code is a compilation scheme heavily based on the \ac{SAM} interpreter. Execution is done sequentially and deterministically, with the compiler generating code that closely resembles the execution of the \ac{SAM}, implementing continuation-based coroutining.
Associated to each process is a block of data called an environment, similar to a stack frame of traditional imperative languages. This environment stores the continuations of each session, as well as any data received through them. Whenever a process finishes writing data to a session, a jump is made to its other end point, which will read the data and continue execution.

At a high-level, the compilation scheme is divided into four phases. First, the source \ac{CLASS} program is parsed and type-checked. This phase is not discussed in this report, as its implementation has been reused from the previous work \cite{Rocha:2023, Caires:2024} on the \ac{SAM} and \ac{CLASS} interpreters. Then, the program is compiled to an intermediate representation (\ac{IR}) which abstracts some details of the target language. Next, the \ac{IR} is analyzed and optimized by a series of passes. Finally, the optimized \ac{IR} is translated to C code, which is then compiled to native code using the \acf{GCC}.

The \ac{IR} optimization phase is left for the next chapter, while this chapter focuses on presenting the base \ac{IR} and compilation scheme. This chapter is divided into five four. In the first section, the \ac{IR}'s constructs and its abstract machine are defined. In the second section, the compilation of \ac{CLASS} to the \ac{IR} is described. Finally, in the third section, an overview is given on how the \ac{IR} is translated to C code.

\section{Intermediate Representation}

In this section, we introduce the intermediate representation (\ac{IR}) used by the compiler and its abstract machine. The \ac{IR} is an imperative stateful language meant to act as a bridge between the high-level session-based language \ac{CLASS} and the low-level imperative language C.

We start by describing the syntax and layout of \ac{IR} programs. Then, we define the data structures and other necessary concepts used by the abstract machine, as we go through the semantics of each \ac{IR} instruction.

A complete specification of the \ac{IR} can be found in \Cref{chapter:appendixIR}.

\subsection{Program Structure}

An \ac{IR} program is a set of named processes, each consisting of a header and a body. The header contains the number of end points of the process, the names of the type parameters, the names of the session parameters, the names of the data sections and their data layouts, and finally, the data to be dropped when the process finishes execution (see \Cref{fig:IRProcessExample} for an example).

\begin{figure}[H]
    \begin{lstlisting}[language=IR]
foo:
    end points: 1
    types: none
    sessions: s0
    data d0: <[]> (s0)
entry:
    printLine("close")
    finishSession(s0, end point)

main:
    end points: 2
    types: none
    sessions: s0
    data d0: <[]> (s0)
entry:
    initializeSession(s0, cut_lhs_0, d0)
    callProcess(foo, end point, s0 <- s0)
cut_lhs_0:
    printLine("wait")
    popTask(end point)
    \end{lstlisting}
    \caption{An example \ac{IR} program with two processes}
    \label{fig:IRProcessExample}
\end{figure}

\subsection{Environments}

Each process invocation creates an \textit{environment} which holds its dynamic state, similarly to a stack frame in a traditional function call. This state includes the number of remaining end points, the properties of any polymorphic types, the state of all sessions, the local data being used by the process and finally, a bitset indicating which data locations should be dropped when the process terminates.

All instructions in a process implicitly operate on its environment. When an instruction refers to a session or a data location, it is referring to the respective section in the current environment.

\subsection{Sessions}

\ac{IR} sessions are similar to both return addresses and coroutines. They are used to model the flow of control between two end points of a \ac{CLASS} channel.
A session is described by four properties: a reference to its opposite end point session (which might be itself), a data location to write to when sending data to it, and a continuation address and environment to switch to when passing control to it.

Sessions are stored in the environment of each process, and are referenced by their index in the environment's session section. For example, \texttt{s0} refers to the first session in the environment, \texttt{s1} to the second session, and so on. Additionally, references to sessions can be passed around as data through the \texttt{session} slot.

\begin{figure}[H]
    \centering
    \begin{lstlisting}[language=IR]
entry:
    initializeSession(s0, rhs, d0) // 1
    continueSession(s0, cont)      // 2
rhs:
    finishSession(s0)              // 3
cont:
    ...                            // 4
    \end{lstlisting}
    \caption{Example flow control using sessions}
    \label{fig:IRSessionFlowControlSimple}
\end{figure}

Program flow through session is manipulated through three main instructions: \texttt{initializeSession}, \texttt{continueSession} and \texttt{finishSession}. An example flow of control through a session is shown in \Cref{fig:IRSessionFlowControlSimple}.

The \texttt{initializeSession(s0, rhs, d0)} instruction initializes the session \texttt{s0} in the current environment such that its remote counterpart (opposite end point) is itself, its writing data location is \texttt{d0} and its continuation is the block \texttt{rhs} on the current process.

Then, the \texttt{continueSession(s0)} instruction is executed, which jumps to the continuation of \texttt{s0} and replaces the continuation of \texttt{s0}'s remote session (itself in this case) with the block \texttt{cont} on the current process. This means that when \texttt{s0} is continued again, control will be passed to \texttt{cont}.

Finally, the \texttt{finishSession(s0)} instruction is executed, which passes control to the continuation of \texttt{s0}, which is \texttt{cont}. Additionally, it detaches \texttt{s0} from its remote session by setting its remote session reference to \texttt{NULL}. This instruction is similar to continue, but with the difference that flow of control won't ever return.

\begin{figure}[H]
    \centering
    \begin{lstlisting}[language=IR]
entry:
    initializeSession(s0, lhs, d0) // 1
    initializeSession(s1, rhs, d1) // 2
    forwardSessions(s0, s1)        // 3
lhs:
    ...                            // 5
rhs:
    finishSession(s1)              // 4
    \end{lstlisting}
    \caption{Example flow control using session forwarding}
    \label{fig:IRSessionFlowControlForwarding}
\end{figure}

Additionally, the instruction \texttt{forwardSessions} is used to model forwarding of \ac{CLASS} channels. An example of its use is shown in \Cref{fig:IRSessionFlowControlForwarding}.

In this example two sessions \texttt{s0} and \texttt{s1} are initialized as before. Then, the \texttt{forwardSessions(s0, s1)} instruction is executed, which makes the remote session of \texttt{s0} be the remote session of \texttt{s1}, and vice-versa. After linking the sessions, the instruction jumps to the continuation of its second argument (\texttt{s1} in this case, and thus, the block \texttt{rhs}).

On the block \texttt{rhs}, the instruction \texttt{finishSession(s1)} is executed, which just like before passes control to the continuation of \texttt{s1}, which was set to \texttt{lhs} by the forwarding instruction.

\subsection{Data Storage}

All data in \ac{IR} programs (excluding meta data such as sessions and type information) is stored in one of two places: within a local data section on an environment, or within a cell. Local data sections are referenced by identifiers such as $d0$ and $d1$, similarly to sessions, while cells are referenced through cell slots (see \Cref{subsection:Cells}).

\subsubsection{Slots} \label{subsection:IRSlots}

\ac{IR} programs are completely platform agnostic, and thus, data types and addresses are described through their composition of \textit{slots}, instead of concrete byte sizes and offsets. A \textit{slot} is a unit of data in the \ac{IR}. There are many types of slots, from primitive types such as $\ints$ and $\strings$ to more domain-specific types such as $\sessions$ and $\exponentials$ (see \Cref{chapter:appendixIR} for a complete list).

In order to model branching data types (such as choices from \ac{CLASS}), data layouts are described through trees of slots, instead of plain sequences. These trees branch on the choices taken (see \Cref{fig:IRSlotTreeExample} for an example).

\begin{figure}[H]
    \centering
    \begin{lstlisting}[language=IR]
// Leaf node, representing no data.
[]

// Integer, followed by a boolean, followed by a string.
[int, bool, string]

// Integer, followed by a tag, followed by a boolean if the tag is 0,
// or by a string if the tag is 1.
[int, tag[[bool] | [string]]]
    \end{lstlisting}
    \caption{Examples of slot trees}
    \label{fig:IRSlotTreeExample}
\end{figure}

Local data sections are described by their identifier and by a sum of slot trees describing the possible layouts of data that can be stored on them. For example, $data\ d0:\ \combst{\seqst{\ints; \bools} + \seqst{\strings}}$ describes a local data section $d0$ which can either store an integer followed by an $\ints$, or a $\strings$. Additionally, data sections may be associated to a session, indicating that a session passed to a process invocation should receive data from that data section (e.g., $data\ d0:\ \combst{\seqst{\ints}}\ (s0)$).

\subsubsection{Addressing Data} \label{subsection:DataAddressing}

A \textit{data location} is defined by a base address and an offset. The base address can be either a local data section (e.g., $d0$), remote data referenced by a session (e.g., $s0$), or a cell location (e.g., $c(d0)$). The offset is defined by two slot trees, written as $\offs{P}{F}$, $P$ describing the data being skipped (named \textit{past}), and $F$ describing the data that may come next (named \textit{future}). The future slot tree is necessary to determine the alignment of the data being accessed. For example, $\offs{\seqst{\ints; \bools}}{\seqst{\strings}}$ is an offset which skips an integer followed by a boolean, and expects a string to come next. When the past slot tree is empty, the future slot tree is unnecessary, and the offset is simply written as $0$. 

Locations are written as $\datal{B}{P}{F}$, where $B$ is the base address, and $\offs{P}{F}$ is the offset. For example, the location $\datal{d0}{\seqst{\ints}}{\seqst{\ints}}$ refers to the second integer stored in the local data section $d0$, while the location $\datal{s0}{\seqst{\bools}}{\seqst{\strings}}$ refers to a string stored in the remote data of the session $s0$, after a boolean. When the offset is zero, it is omitted for brevity. For example, $d0$ as a data location is equivalent to $\datal{d0}{\nilst}{\nilst}$ or just $d0.0$ - the future slot trees are unnecessary for empty past slot trees.

\subsection{Moving data}

Sometimes it is necessary to simply move data from one location to another. This is done through the \texttt{moveSlots} instruction, which takes as arguments a destination location, a source location and a slot tree describing the data to be moved. To exemplify, \texttt{moveSlots(s1, $\datal{d0}{\seqst{\bools}}{\seqst{\strings; \ints}}, \seqst{\strings, \ints}$)} moves a string followed by an integer from the local data section \texttt{d0} after a boolean to the start of the remote data section of session \texttt{s1}.

The \texttt{cloneSlots} instruction is used to clone data from one location to another, taking the same arguments. This instruction is only defined for cloneable slots, which are $\ints$, $\bools$, $\strings$, $\tags$ and $\exponentials$. In addition to copying the memory, it also increments reference counts for exponentials and allocates new strings as necessary. It used to model the exponential part of \ac{CLASS}.

\subsection{Dropping data} \label{subsection:IRDroppingData}

Although \ac{CLASS} is a linear language, it also supports both affine and exponential types. Thus, it is necessary to be able to drop some types of data if they are no longer needed. To this end, two instructions are provided: \texttt{dropSlots} and \texttt{deferDrop}.

The \texttt{dropSlots} instruction takes as arguments a data location and a slot tree, and immediately drops the passed data. This instruction is only defined for droppable slots, which are $\ints$, $\bools$, $\strings$, $\tags$, $\sessions$, $\cells$ and $\exponentials$. 

When a $\sessions$ slot is dropped, the session is expected to follow the affine protocol, defined is as follows: when the session is to be deallocated, a tag with value $0$ is written to indicate the session will not be used, and control is passed to it with a continue instruction. Then, the session should finish, passing control back to the cell deallocator.

The \texttt{deferDrop} instruction takes as argument a drop bit identifier, and sets the respective bit. Drop bits are defined in the process header, and indicate a data location and a slot tree to be dropped when the process terminates, if their respective bit is set. Drop bit identifiers start with \texttt{e} followed by a number, such as \texttt{e0} and \texttt{e1}.

\begin{figure}[H]
    \centering
    \begin{lstlisting}[language=IR]
main:
    data d0: <[string; string]>
    drop e0: [string] at d0.[string]~[string]
entry:
    ... // something that initializes both strings at d0
    deferDrop(e0) // the second string will be dropped later
    dropSlots(d0, [string]) // the first string is dropped now
    \end{lstlisting}
    \caption{Dropping data directly and through a drop bit}
\end{figure}

\subsection{Expressions}

Expressions are tree-like structure which compute basic values (\texttt{int}, \texttt{bool} or \texttt{string}) without side effects. Examples of expressions include arithmetic operations, comparisons, boolean operations, string concatenation and data access. A full specification of expressions can be found in \Cref{chapter:appendixIR}.

While expressions do not occur as instructions, they are used as arguments to some instructions, such as \texttt{writeExpression}, \texttt{branchExpression} and \texttt{print}.

An example of an expression is \texttt{(move(d0, int) + 3) + clone(d1, string)}. This expression reads an integer from \texttt{0}, adds \texttt{3} to it, then reads a cloned string from \texttt{d1} and finally concatenates the two. The \texttt{move} and \texttt{clone} expressions read data from the respective locations, with the difference that \texttt{move} consumes the data (cannot be read again) while \texttt{clone} leaves the data intact.

In the particular case of integers and booleans, both \texttt{move} and \texttt{clone} are equivalent, as these types are trivially copyable. For strings, however, the difference is important, as \texttt{clone} allocates a new string and copies the contents, while \texttt{move} simply uses the existing allocation.

\subsection{Input and output}

The \texttt{print} and \texttt{printLine} instructions are used to print values to the standard output. They take a single argument expression, which is evaluated and printed. The \texttt{printLine} instruction additionally appends a new line after printing the value.

The \texttt{writeScan} instruction is used to read values from the standard input. It takes as arguments a data location to store the read value to and a slot describing the type of data to be read.

\begin{figure}[H]
    \centering
    \begin{lstlisting}[language=IR]
main:
    data d0: <[int; int]>
entry:
    writeScan(d0, int)
    writeScan(d0.[int]~[int], int)
    printLine(move(d0, int) + move(d0.[int]~[int], int))
    ...
    \end{lstlisting}
    \caption{Reading two integers from standard input and printing their sum}
\end{figure}

\subsection{Writing and binding sessions}

\begin{sloppypar}
References to sessions (\texttt{session} slots) can be written to and read from data locations using the \texttt{writeSession} and \texttt{bindSession} instructions, respectively. The first instruction takes a data location and a session identifier as arguments. The second takes a data location to read from, a session identifier to bind to, and a data location to set as the writing location of the session's remote end point.
\end{sloppypar}

An example of writing a session to local data and then binding it to another session is shown in \Cref{fig:IRSessionDataExample}. First, a session \texttt{s1} is initialized with its remote end point being itself, its writing location being \texttt{d1} and its continuation being the block \texttt{s1\_cont}. Then, a reference to the remote counterpart of session \texttt{s1} is written to the local data section \texttt{d0}. Next, the reference stored at \texttt{d0} is read and bound to the session \texttt{s2}, with its remote writing location being modified to \texttt{d2}. At this point, the remote session of \texttt{s2} is \texttt{s1}, and vice-versa. Finally, \texttt{s2} is continued, jumping to \texttt{s1\_cont}, and setting the continuation of \texttt{s1} to the current block. Thus, when executing \texttt{finishSession(s1)}, control will jump to \texttt{s2\_cont}.

\begin{figure}[H]
    \centering
    \begin{lstlisting}[language=IR]
    initializeSession(s1, d1, s1_cont) // 1
    writeSession(d0, s1)               // 2
    bindSession(d0, s2, d2)            // 3
    continueSession(s2, s2_cont)       // 4
s1_cont:
    finishSession(s1)                  // 5
s2_cont:
    ...                                // 6
    \end{lstlisting}
    \caption{Writing a session to local data and binding it to another session}
    \label{fig:IRSessionDataExample}
\end{figure}

\subsection{Tags and branching} \label{subsection:Branching}

Branching is done through the \texttt{branchTag} and \texttt{branchExpression} instructions. The first instruction branches on a tag stored a data location, similarly to a switch statement in C, and the second instruction branches on the result a boolean expression, similarly to an if statement in C.

Each possible branch is defined by a code location (block label) and an integer indicating how many end points are behind that branch. This is used to keep track of the number of remaining end points in the environment, and is explained in more detail in \Cref{subsection:EndPoints}.

Tags are written to data locations through the \texttt{writeTag} instruction, which takes as arguments a data location and a constant tag value.

\begin{figure}[H]
    \centering
    \begin{lstlisting}[language=IR]
main:
    data d0: <tag[[int] | [string]]>
entry:
    writeTag(d0, 1)
    ...
    branchTag(d0, if_int:1, if_string:1)
if_int:
    ...
if_string:
    ...
    \end{lstlisting}
    \caption{Branching on a tag}
\end{figure}

\subsection{Tasks}

Tasks are a flow control mechanism completely separate from sessions, used to model the mix construct from \ac{CLASS}. For each thread of execution, there's a stack of tasks, where each task stores a continuation to be executed when the current task is finished.

Tasks can be pushed to the stack through the \texttt{pushTask} instruction, which takes a code location as argument. To finish the current task and jump to the continuation of the topmost task, the \texttt{popTask} instruction is used.

At the start of a thread, the task stack is initialized with a single task responsible for terminating the thread. Thus, the last instruction to be executed in a thread is always \texttt{popTask}.

\subsection{Jumps}

A simple jump to a constant code location is also possible through the \texttt{jump} instruction. This instruction is useful to merge control flow after branching, for example.

\subsection{Cells} \label{subsection:Cells}

Cells are used to model the shared state part of \ac{CLASS}. They're mutable memory locations which can be safely shared between multiple threads of execution. Cells are referenced through \texttt{cell} slots, and maintain a mutex to ensure safe concurrent access, as well as a reference count to guarantee proper deallocation.

A new cell can be created through the \texttt{writeCell} instruction, which takes as arguments a data location to write the \texttt{cell} slot to, and a slot tree describing the data that will be stored in the cell. The reference count of a cell can be incremented and decremented through the \texttt{acquireCell} and \texttt{releaseCell} instructions, respectively. Finally, the mutex of a cell can be locked and unlocked through the \texttt{lockCell} and \texttt{unlockCell} instructions, respectively.

The data stored in a cell can be accessed through any data location, by using a cell location as the base address, as explained in \Cref{subsection:DataAddressing}.

\begin{figure}[H]
    \centering
    \begin{lstlisting}[language=IR]
main:
    data d0: <cell[int]>
entry:
    writeCell(d0, [int]) // create a cell with space for an integer
    lockCell(c(d0)) // lock the cell's mutex
    writeExpression(c(d0), 42) // write 42 to the cell
    unlockCell(c(d0)) // unlock the cell's mutex
    acquireCell(c(d0)) // increase reference count to 2
    releaseCell(c(d0)) // decrease reference count to 1
    releaseCell(c(d0)) // decrease reference count to 0, deallocating the cell
    \end{lstlisting}
    \caption{Creating and using a cell}
\end{figure}

\subsubsection{Cell Slot Arguments} \label{subsubsection:CellSlotArguments}

While other slots are represented solely by their type, \texttt{cell} require a slot tree argument. This slot tree describes the data that will be stored in the cell, and is necessary to properly drop the data when the cell is deallocated. This might involve invoking the affine protocol on sessions, deallocating strings, or even dropping other cells.

\subsection{Threads} \label{subsection:IRThreads}

New threads of execution can be launched through the \texttt{launchThread} instruction, which takes as arguments a code location to start executing at. The new thread starts with a single task responsible for terminating it, and shares the same environment as the parent thread, but a separate task stack.

Accesses to the environment are not synchronized except for reference and end point counting, and it is the responsibility of the \ac{IR} generator to ensure different threads do not access the same data concurrently.

\subsection{Polymorphism}

Processes in the \ac{IR} may be parametrized by polymorphic types. These types are identified by names which start with \texttt{t} followed by a number, such as \texttt{t0} and \texttt{t1}. Types may be used directly as slots in slot trees (e.g.: $\seqst{t_2; \ints; t_0}$).

Types can be written to data locations through the \texttt{writeType} instruction, producing a \texttt{type} slot. This slot can then be read by a process invocation instruction, such as \texttt{callProcess}. The \texttt{writeType} instructions takes as arguments a data location to write to and a slot tree describing the data layout of the type.

For example, \texttt{writeType(d0, [int; t1])} writes a type to \texttt{d0} which stores first an integer, and then data of type \texttt{t1}, which is a polymorphic type in the current environment.

\subsection{Process calls}

Processes can be invoked through the \texttt{callProcess} instruction, which takes as arguments the name of the process to invoke (might be the same process), and a list of type, session and data arguments to initialize the new process environment.

A type argument is described by the target type identifier on the new process, and by either a data location containing a \texttt{type} slot, or by a slot tree describing the type directly. For example, \texttt{t0 <- [int]} initializes the type parameter \texttt{t0} of the new process to be a type which stores an integer, while \texttt{t1 <- d3.[int]~[type]} initializes \texttt{t1} to the type stored after an integer at \texttt{d3}.

A session argument is described by the target session identifier on the new process, by either a source session identifier or a data location containing a \texttt{session} slot, and finally, by a data offset to apply to the writing location of the session. For example, \texttt{s0 <- s1+[int]~[int]} initializes the session parameter \texttt{s0} of the new process to be the same as session \texttt{s1} of the current process, but with its writing location pointing at $\datal{s1}{\seqst{\ints}}{\seqst{\ints}}$. If the offset is zero, it is omitted for brevity.

Finally, a data argument is described by the target data identifier on the new process, by a source data location, by a slot tree describing the data, and by a flag indicating whether to clone the data or not. For example, \texttt{d0 =[string] d1} writes to \texttt{d0} a clone of the string stored at \texttt{d1}, while \texttt{d1 <-[string] d2} moves to \texttt{d1} the string stored at \texttt{d2} without cloning it.

\begin{figure}[H]
    \centering
    \begin{lstlisting}[language=IR]
foo:
    types: t0
    sessions: s0
    data d0: <[t0]>
entry:
    moveSlots(s0, d0, [t0]) // 4
    finishSession(s0)       // 5

main:
    sessions: s0
    data d0: <[int]>
entry:
    initializeSession(s0, d0, rhs)                                     // 1
    writeExpression(s0, 42)                                            // 2
    callProcess(foo, t0 <- [int], s0 <- s0, d0 <-[int] d0) // 3
rhs:
    printLine(move(d0, int)) // 6
    ...
    \end{lstlisting}
    \caption{Calling a process with type, session and data arguments}
    \label{fig:IRProcessCallExample}
\end{figure}

An example of calling a process with type, session and data arguments is shown in \Cref{fig:IRProcessCallExample}. First, the session \texttt{s0} is initialized to have its remote end point be itself, its writing location be \texttt{d0} and its continuation be the block \texttt{rhs}. Then, the integer \texttt{42} is written to the remote data of \texttt{s0} (which is \texttt{d0}) Next, the process \texttt{foo} is called, with its type parameter \texttt{t0} being set to an integer, its session \texttt{s0} being set to the current session \texttt{s0}, and its data section \texttt{d0} being initialized by a move of the local data section \texttt{d0}. Inside \texttt{foo}, the integer stored at its local data section \texttt{d0} (received from \texttt{main}) is moved to the remote data of its session \texttt{s0} (which is still \texttt{d0} from the \texttt{main} process), and then, the session is finished, passing control back to the block \texttt{rhs} of the calling process. Finally, the integer stored at \texttt{d0} is printed.

\subsection{Exponentials}

Exponentials in \ac{CLASS} can be modeled through the \texttt{exponential} slot, which stores the necessary data to dynamically call a process closure, taking care of any captured data. Exponentials are created through the \texttt{writeExponential} instruction, which takes the same arguments as \texttt{callProcess}, with the exception that session arguments are not allowed, and that passed data must be cloneable. Additionally, the referenced process is expected to follow the exponential process convention.

The convention for exponential processes is simply that they must take a single session argument, \texttt{s0}, which will be used to communicate with the caller. This session is initialized by the caller when the \texttt{callExponential} instruction is called. This instruction takes as arguments a data location containing an \texttt{exponential} slot, a session identifier to bind the exponential process session to, and a data location to set as the writing location on the exponential process. See \Cref{fig:IRExponentialCallExample} for an example of creating and calling an exponential.

One important detail is that exponentials are reference counted data, and the \texttt{exponential} slot must be dropped when no longer needed, or at the end of the process, if it was not moved. This is done through the \texttt{dropSlots} or \texttt{deferDrop} instructions, as explained in the \Cref{subsection:IRDroppingData} section. Drops of exponentials decrement this reference count, and clones increment it.

\begin{figure}[H]
    \centering
    \begin{lstlisting}[language=IR]
foo:
    types: t0
    sessions: s0
    data d0: <[t0]>
entry:
    moveSlots(s0, d0, [t0]) // 6, 10
    finishSession(s0)       // 7, 11

main:
    sessions: s0 s1
    data d0: <[string]>
    data d1: <[exponential]>
    data d2: <[string]>
    data d3: <[string]>
    drop e0: [exponential] at d1
entry:
    writeExpression(d0, "Hello, World!")                                    // 1
    writeExponential(d1, foo, t0 <- [string], d0 <-[string] d0) // 2
    deferDrop(e0)               // 3
    callExponential(d1, s0, d2) // 4
    callExponential(d1, s1, d3) // 4
    continueSession(s0, rhs0)   // 5
rhs0:
    printLine(move(d2, string)) // 8
    continueSession(s1, rhs1)   // 9
rhs1:
    printLine(move(d3, string)) // 12
    ...
    \end{lstlisting}
    \caption{Calling an exponential with type and data arguments}
    \label{fig:IRExponentialCallExample}
\end{figure}

In \Cref{fig:IRExponentialCallExample}, the process \texttt{foo} is a simple polymorphic process which takes a type parameter \texttt{t0}, a session \texttt{s0} and a data section \texttt{d0}. It moves data of type $\seqst{\strings}$ from its local data section \texttt{d0} to the remote data of its session \texttt{s0}, and then finishes the session, passing control back to the caller. When called as an exponential, \texttt{s0} will be initialized by the caller.

The \texttt{main} process first stores a new string at \texttt{d0}, and then creates a new exponential at \texttt{d1} for the process \texttt{foo}, moving the created string to it. This string will be cloned every time the exponential is called, and dropped when the exponential is freed. The exponential will be dropped at the end of the process through the drop bit \texttt{e0}, which is set by the \texttt{deferDrop(e0)} instruction.

The exponential is then called twice, first binding its session to \texttt{s0} and its writing location to \texttt{d2}, and then binding its session to \texttt{s1} and its writing location to \texttt{d3}. Each call will clone the string stored at \texttt{d0} and pass it to the exponential process. Control is not immediately passed to the exponential process. Instead, the caller only passes control to the exponential process when it continues either \texttt{s0} or \texttt{s1}. This is done through the \texttt{continueSession} instructions, which jump to the continuation of the respective session, which was set to the entry block of the exponential process by the \texttt{callExponential} instructions.

Finally, within the exponential process, the string is moved from the local data section \texttt{d0} to the remote data of the session \texttt{s0} (which is either \texttt{d2} or \texttt{d3} of the caller, depending on which call was continued), and then, the session is finished, passing control back to the caller. The caller then prints the received string.

\subsection{End points} \label{subsection:EndPoints}

One important aspect of the \ac{IR} which hasn't been explained yet is the concept of end points. These are used to ensure proper deallocation of the process environment when there's no more code left to execute within it. Each environment maintains a count of the number of remaining end points, which is specified in the process header (omitted in previous examples for brevity).

Each instruction which passes control to another block without a continuation might be marked as an end point. These instructions are \texttt{finishSession}, \texttt{forwardSessions}, \texttt{popTask} and \texttt{callProcess}. When such an instruction is executed, the environment's end point count is decremented. When it reaches zero, the environment is deallocated, and any remaining data is dropped according to the drop bits defined in the process header. A simple example of this mechanism is shown in \Cref{fig:IREndPointsExample1}.

\begin{figure}[H]
    \centering
    \begin{lstlisting}[language=IR]
main:
    end points: 2
    sessions: s0
    data d0: <[]>
entry:
    initializeSession(s0, d0, rhs)
    finishSession(s0, end point) // end points decremented to 1
rhs:
    popTask(end point) // end points decremented to 0, environment is freed
    \end{lstlisting}
    \caption{Decrementing the end point count and deallocating the environment}
    \label{fig:IREndPointsExample1}
\end{figure}

The only other instructions which interact with the end point mechanism are the branching instructions, as mentioned in \Cref{subsection:Branching}. Each possible branch of a branching instruction is marked with the number of end points that will be behind that branch. When a given branch is taken, the environment's end point count is decremented by the difference between the maximum end point count across all branches and by end point count of the chosen branch. An example of this mechanism is shown in \Cref{fig:IREndPointsExample2}.

\begin{figure}[H]
    \centering
    \begin{lstlisting}[language=IR]
foo:
    end points: 2
    data d0: <[int]>
entry:
    branchExpression(move(d0, int) > 0, if_true:2, if_false:1) // max = 2
if_true:
    ... // end points remains 2 (max - curr = 2 - 2 = 0)
if_false:
    ... // end points decremented by 1 (max - curr = 2 - 1 = 1)
    \end{lstlisting}
    \caption{Modifying the end point count through branching}
    \label{fig:IREndPointsExample2}
\end{figure}

\section{Generating IR programs from CLASS}

This section formally defines the base (unoptimized) translation from \ac{CLASS} to the \ac{IR} defined in the previous section. It starts defining notation for representing and updating the state of the compiler. Then, it defines how session types are mapped to the slot trees introduced in \Cref{subsection:IRSlots}. Next, an overview of the instruction generation scheme is presented, and some important concepts are defined. Finally, the translation from \ac{CLASS} processes to \ac{IR} programs is incrementally defined, starting from simple processes and adding more constructs until the whole language is covered.

This compilation scheme assumes well-typed \ac{CLASS} processes are given as input, along with the required type information of free names in all processes. The notation $x: A$ is used to indicate that the type checker has inferred that channel $x$ is a free name of type $A$ on some process.

\subsection{Compiler state}

To compile \ac{CLASS} processes to \ac{IR} programs, the compiler must keep track of information regarding the state of each type and channel name within the process being compiled. This state is referred to as the environment, and is represented by a tuple $(\theta, \Delta)$, where $\theta$ is the type environment and $\Delta$ is the channel environment.

\subsubsection{Type environments} \label{subsubsection:TypeEnvironments}

A type environment $\theta$ is a mapping from type variables to a tuple $(p, t)$ where $p$ is the polarity of the type variable ($+$ or $-$) and $t$ is its identifier (e.g, $t_0$, $t_1$) within the header of the process being compiled. $\polarityof{\theta}{X}$ and $\typeidof{\theta}{X}$ are defined as the polarity and identifier of type variable $X$ in the type environment $\theta$, respectively.

A type with positive polarity is one for which sessions of that type are used to write data, while a type with negative polarity is one for which sessions of that type are used to read data. The function $\polarityof{\theta}{A}$ is further extended to all session types $A$ through pattern matching, as shown in \Cref{fig:TypePolarity}. For the base case of type variables, the polarity is looked up in the type environment.

\begin{figure}
    \centering
    \begin{align*}
        &\polarityof{\theta}{\lintt} = + && \polarityof{\theta}{\colintt} = - && \polarityof{\theta}{\lboolt} = + && \polarityof{\theta}{\colboolt} = - \\
        &\polarityof{\theta}{\lstringt} = + && \polarityof{\theta}{\colstringt} = - && \polarityof{\theta}{\closet} = + && \polarityof{\theta}{\waitt} = - \\
        &\polarityof{\theta}{\choicet{\dots}} = + && \polarityof{\theta}{\offert{\dots}} = - && \polarityof{\theta}{\sendt{A}{B}} = + && \polarityof{\theta}{\recvt{A}{B}} = - \\
        &\polarityof{\theta}{\bangt{A}} = + && \polarityof{\theta}{\whynott{A}} = - && \polarityof{\theta}{\sendtyt{X}{A}} = + && \polarityof{\theta}{\recvtyt{X}{A}} = - \\
        &\polarityof{\theta}{\affinet{A}} = + && \polarityof{\theta}{\coaffinet{A}} = - && \polarityof{\theta}{\statet{A}} = + && \polarityof{\theta}{\usaget{A}} = - \\
        &\polarityof{\theta}{\rect{X}{A}} = + && \polarityof{\theta}{\corect{X}{A}} = - && \polarityof{\theta}{\statelt{A}} = + && \polarityof{\theta}{\usagelt{A}} = - \\
    \end{align*}
    \caption{Polarity of session types}
    \label{fig:TypePolarity}
\end{figure}

\subsubsection{Channel environments}

A channel environment $\Delta$ is a mapping from channel names to channel states. Each channel state is a tuple $(s, d, o)$, where $s$ is the session identifier associated with the channel (or $0$ if there's no session), $d$ is the identifier of the local data section from which data is read (or $0$ if there's no data section), and $o$ is the data offset to apply to the location of the next write or read operation on the channel. The data offset is represented as a slot tree, as explained in \Cref{subsection:DataAddressing}.

Given a channel environment $\Delta$, $\Delta[x \mapsto (s, d, o)]$ is defined as the channel environment $\Delta$ where channel $x$ is mapped to the channel state $(s, d, o)$, overwriting any previous mapping for $x$. Additionally $\sessionof{\Delta}{x}$, $\localdataof{\Delta}{x}$ and $\offsetof{\Delta}{x}$ are defined as the session identifier, local data identifier and data offset associated with channel $x$ in the channel environment $\Delta$, respectively.

\subsubsection{Writing and reading locations}

For convenience, both $\readlocof{\Delta}{x}$ and $\writelocof{\Delta}{x}$ are defined. The first provides the current data location from which data of channel $x$ is to be read, taking into account the offset, and is defined as $\readlocof{\Delta}{x} = \dataol{\localdataof{\Delta}{x}}{\offsetof{\Delta}{x}}$. The second is the current data location into which data for channel $x$ is to be written, and is defined as $\writelocof{\Delta}{x} = \dataol{\localdataof{\Delta}{x}}{\offsetof{\Delta}{x}}$.

\subsubsection{Advancing the offset} \label{subsubsection:AdvancingOffsets}

When more than one slot is read from or written to a channel in succession, the channel's offset must be advanced accordingly. To this end, first the concatenation of offsets must be defined. Given two offsets $o_1 = (P_1, F_1)$ and $o_2 = (P_2, F_2)$, their concatenation is defined as $\concatst{o_1}{o_2} = (\concatst{P_1}{P_2}, F_2)$, where $\concatst{P_1}{P_2}$ is the concatenation of the two slot trees. In other words, when two offsets are concatenated, the resulting offset has concatenates both of the past slot trees, and keeps the future slot tree of the second offset. Slot tree concatenation is defined as replacing every leaf of the first slot tree with the second slot tree.

To simplify the definitions in the next sections, $\advancech{\Delta}{x}{p}{\theta}{S}{A}$ is defined as the channel environment $\Delta$ where the offset of channel $x$ is either advanced by a given offset, or reset. The decision is made based on whether the polarity $p$ (either $+$ or $-$) matches the polarity of the type $A$ in the type environment $\theta$. If they match, the new offset is computed as $\concatst{\offsetof{\Delta}{x}}{(\slotoffset{S}{\slots{\theta}{A}})}$, where $S$ is a slot tree describing the data which was just read or written, and $\slots{\theta}{A}$ is the slot tree corresponding to type $A$, which will be defined in the next section. Otherwise the new offset is $0$. Formally:

$$
\advancech{\Delta}{x}{p}{\theta}{S}{A} = \begin{cases}
    \Delta[x \mapsto (\sessionof{\Delta}{x}, \localdataof{\Delta}{x}, \concatst{\offsetof{\Delta}{x}}{(\slotoffset{S}{\slots{\theta}{A}})})] & \text{if } p = \polarityof{\theta}{A} \\
    \Delta[x \mapsto (\sessionof{\Delta}{x}, \localdataof{\Delta}{x}, 0)] & \text{otherwise}
\end{cases}
$$

\subsection{Mapping session types to slot trees}

To translate \ac{CLASS} processes to \ac{IR} programs, we need to define which data is sent and received over sessions. As explained in \Cref{subsection:IRSlots}, in the \ac{IR} data types are represented through slot trees. Thus, we need to define how session types are mapped to slot trees.

We start by defining this mapping for session types $A$ which are strictly positive. Let $\slots{\theta}{A}$ be the slot tree corresponding to the data exchanged over a session type $A$ within type environment $\theta$. This function is defined through pattern matching on session types in \Cref{fig:SessionTypeToSlotTree}.

\begin{figure}[H]
    \centering
    \begin{align*}
        &\slots{\theta}{\closet} = \nilst && \slots{\theta}{\lintt} = \seqst{\ints} \\
        &\slots{\theta}{\lboolt} = \seqst{\bools} && \slots{\theta}{\lstringt} = \seqst{\strings} \\
        &\slots{\theta}{\sendt{A}{B}} = \seqst{\sessions; \slots{\theta}{B}} && \slots{\theta}{\bangt{A}} = \seqst{\exponentials} \\
        &\slots{\theta}{\sendtyt{X}{A}} = \seqst{\types; \sessions; \tags} && \slots{\theta}{\affinet A} = \seqst{\sessions} \\
        &\slots{\theta}{\statet A} = \seqst{\cells} && \slots{\theta}{\statelt A} = \seqst{\cells} \\
        &\slots{\theta}{\choicet{\labelt a A\ (\dots)\ \labelt z Z}} = \tagst{\slots{\theta}{A} \texttt{|}\ \dots\ \texttt{|}\slots{\theta}{Z}} && \slots{\theta}{X} = \seqst{\typeidof{\theta}{X}} \text { s.t. } \polarityof{\theta}{X} = + \\
    \end{align*}
    \caption{Mapping from value session types to slot trees}
    \label{fig:SessionTypeToSlotTree}
\end{figure}

To exemplify, consider the session type $\sendt{\lintt}{\sendt{\closet}{\choicet{\labelt a \lstringt\ \labelt b \lboolt}}}$. Its slot tree is $\seqst{\sessions; \sessions; \tagst{\seqst{\strings}\tagsepst \seqst{\bools}}}$, as first a session is sent, then another session is sent, and finally, a tag is sent followed by either a string or a boolean, depending on the sent tag.

To extend this concept to all session types, we must consider that session types may have multiple segments of data exchange, separated by changes in polarity. We can separate a session type into its segments by replacing every continuation of a different polarity with the unit type of that polarity ($\closet$ or $\waitt$) on the current segment and starting a new segment with the continuation.

For example, we can split $\sendt{\lintt}{\recvt{\lstringt}{\closet}}$ into three segments: first $\sendt{\lintt}{\closet}$, then $\recvt{\lstringt}{\waitt}$ and finally  $\closet$. Each segment is either strictly positive or strictly negative. Strictly negative segments can be turned into strictly positive segments by computing their dual types. Thus, we can map each segment to a slot tree through the $\slots{\theta}{A}$ function defined above.

The slot tree corresponding to the first segment of a session type $A$ is designated as its \textit{active slot tree}  (denoted simply by $\slots{\theta}{A}$), and the slot trees corresponding to the remaining segments as its \textit{remainder slot trees}. Additionally, we denote positive segments as \textit{remote slot trees} and negative segments as \textit{local slot trees}, as positive segments correspond to data which is sent (and thus stored remotely), while negative segments correspond to data which is received (and thus stored locally). The local set of slot trees of a session type $A$ is denoted by $\alllocalslots{\theta}{A}$, and the total set of slot trees by $\allslots{\theta}{A}$. Sets of slot trees are termed \textit{layouts}, as they describe the full set of data that may be exchanged over a session.

\subsection{Instruction generation}

The \ac{IR} instructions generated from a \ac{CLASS} process $P$ on a given environment $(\theta, \Delta)$ are denoted by $\getinstructions{\theta}{\Delta}{P}$. This function produces a sequence of \ac{IR}, optionally followed by blocks with more instructions, and is defined incrementally in the following sections.

The generated code implements the write-to-read adjustment of the \ac{SAM}, by placing the instruction \texttt{continueSession} whenever the polarity of a session changes from positive to negative. To simplify the definition of the generation rules, a pseudo-instruction $\texttt{continueIfNegative}(\theta, x, A)$ is used, which is replaced by a \texttt{continueSession} instruction and a new block for its proceeding instructions if and only if $\polarityof{\theta}{A} = -$.

All writing operations to a channel $x$ write to an offset within the remote data of the session, i.e.,  $\writelocof{\Delta}{x}$. All reading operations read from an offset within the local data of the session, i.e., $\readlocof{\Delta}{x}$. Through the mechanism above, whenever the writing side finishes writing data,, the session is continued, and the reading side reads from the same location the writing side wrote to.

For example, consider two processes communicating over a session of type $\sendt{A}{\recvt{B}{\waitt}}$, where the positive channel is named $x$ and the negative channel is named $y$. The writing process, with a channel environment $\Delta^1$ first writes a session of type $A$ to $\writelocof{\Delta^1}{x}$, which is equal to $\readlocof{\Delta^2}{y}$. Since $\recvt{B}{\waitt}$ is negative, a write-to-read adjustment is performed, and thus, a \texttt{continueSession} instruction is executed on $\sessionof{\Delta^1}{x}$.

The reading process, with a channel environment $\Delta^2$, reads a session of type $A$ from $\readlocof{\Delta^2}{y}$. Then, it will write to $\writelocof{\Delta^2}{y}$ a session of type $B$. The writing offset is zero, since we performed a read-to-write adjustment after reading $A$. Finally, it finishes the session $\sessionof{\Delta^2}{y}$, passing control back to the writing process.

The writing process now reads from $\readlocof{\Delta^1}{x}$ a session of type $B$, which is equal to $\writelocof{\Delta^2}{y}$, and the communication is complete.

\subsection{Expression generation}

Apart from generating instructions for processes, \ac{IR} expressions must be generated from \ac{CLASS} expressions when generating instructions for processes such as $\printp$, $\letp$ and $\ifp$. To achieve this, the function $\getexpression{\Delta}{E}$ is defined, which takes a \ac{CLASS} expression $E$ and a channel environment $\Delta$, and produces equivalent an \ac{IR} expression.

This function is defined through pattern matching on \ac{CLASS} expression, as shown below.

\begin{align*}
    &\getexpression{\Delta}{E_1 + E_2} = \getexpression{\Delta}{E_1} + \getexpression{\Delta}{E_2} & & |\ & & \getexpression{\Delta}{E_1 - E_2} = \getexpression{\Delta}{E_1} - \getexpression{\Delta}{E_2} \\
    &\getexpression{\Delta}{E_1 * E_2} = \getexpression{\Delta}{E_1} * \getexpression{\Delta}{E_2} & & |\ & & \getexpression{\Delta}{E_1 / E_2} = \getexpression{\Delta}{E_1} / \getexpression{\Delta}{E_2} \\
    &\getexpression{\Delta}{E_1 \% E_2} = \getexpression{\Delta}{E_1}\,\%\,\getexpression{\Delta}{E_2} & & |\ & & \getexpression{\Delta}{-E} = -\getexpression{\Delta}{E} \\
    &\getexpression{\Delta}{E_1\,\texttt{and}\,E_2} = \getexpression{\Delta}{E_1}\ \texttt{and}\ \getexpression{\Delta}{E_2} & & |\ & & \getexpression{\Delta}{E_1\,\texttt{or}\,E_2} = \getexpression{\Delta}{E_1}\ \texttt{or}\ \getexpression{\Delta}{E_2} \\
    &\getexpression{\Delta}{\texttt{not}\,E} = \texttt{not}\ \getexpression{\Delta}{E} & & |\ & & \getexpression{\Delta}{E_1 < E_2} = \getexpression{\Delta}{E_1} < \getexpression{\Delta}{E_2} \\
    &\getexpression{\Delta}{E_1 \leq E_2} = \getexpression{\Delta}{E_1} \leq \getexpression{\Delta}{E_2} & & |\ & & \getexpression{\Delta}{E_1 > E_2} = \getexpression{\Delta}{E_1} > \getexpression{\Delta}{E_2} \\
    &\getexpression{\Delta}{E_1 \geq E_2} = \getexpression{\Delta}{E_1} \geq \getexpression{\Delta}{E_2} & & |\ & & \getexpression{\Delta}{E_1 == E_2} = \getexpression{\Delta}{E_1} == \getexpression{\Delta}{E_2} \\
    &\getexpression{\Delta}{n} = n \text{ where } n \text{ is an integer literal} & & |\ & & \getexpression{\Delta}{b} = b \text{ where } b \text{ is a boolean literal} \\
    &\getexpression{\Delta}{s} = s \text{ where } s \text{ is a string literal} & & |\ & & \getexpression{\Delta}{x} = \texttt{move}(\readlocof{\Delta}{x}) \text{ where } x \text{ is linear}\\
    &\getexpression{\Delta}{x} = \texttt{clone}(\readlocof{\Delta}{x}) \text{ where } x \text{ is exponential}
\end{align*}

\subsection{End point counting}

One other important aspect of the compilation is counting the number of end points of the generated \ac{IR} code for a given \ac{CLASS} process. This is necessary to ensure that the generated code properly deallocates the process environment when there's no more code left to execute within it, as explained in \Cref{subsection:EndPoints}.

The number of end points in a process $P$ is denoted by $\getendpoints{\theta}{P}$, which is defined by pattern matching on processes. It will be incrementally defined along with the instruction generation rules in the following sections.

\subsection{Process Definitions} \label{subsection:ProcessDefinitions}

A \ac{CLASS} program is a set of process definitions, each with a unique name. For each process definition, a set of resulting \ac{IR} processes is generated, one for each possible combination of polarities of its type parameters. This is necessary to ensure that the polarities of type parameters are always known at compile time, as they greatly affect the flow of control of the generated code.

For each resulting \ac{IR} process, an environment $(\theta, \Delta)$ is initialized. The type environment $\theta$ maps each type parameter name to a tuple $(p, t)$ such that $p$ is the polarity of the type parameter in the current combination, and $t$ is a unique identifier for the type parameter, which is inserted into the process header.

The channel environment $\Delta$ maps each argument channel name to a channel state $(s, d, 0)$. If the argument is linear, $s$ is a fresh session identifier, inserted into the process header. Otherwise, $s = 0$, as exponential channels don't have an associated session. In both cases, $d$ is a fresh local data identifier inserted into the process header, with a layout given by $\alllocalslots{\theta}{A}$, where $A$ is the type of the argument channel. This reserves space for all the data which will be received over the channel. The offset is initialized to $0$, as no data has been read or written yet.

For each of the exponential arguments of the process, an always set drop bit is also added to the process header, to ensure that the exponential data is dropped when the process ends.

Finally, we generate the instructions for the process body $P$ (starting by the entry block) through $\getinstructions{\theta}{\Delta}{P}$, and set the number of end points of the process to $\getendpoints{\theta}{P}$. As instructions are generated, the process header is modified incrementally to include any new sessions, local data sections or drop bits which are necessary for the generated code to work.

\subsection{Cut, Parallel and Empty}

We start by defining the compilation of a $\cutp$ process, which can be seen below. We assume, without loss of generality, that $A$ is a positive type, and thus, that $P$ will be writing to $x$ before $Q$ executes - if $A$ were negative, we could simply swap $P$ and $Q$. The number of end points of the generated code is the sum of the number of end points of both $P$ and $Q$.

First, a fresh session identifier $s$ is generated and added to the process header. Additionally, a data section with unique identifier $d$ is added to the header, with a layout given by $\allslots{\theta}{A}$. In other words, $d$ reserves space for all data which will be both sent and received over the session of type $A$. This is necessary as both sides of the cut share the same data section.

The process initializes the session $s$ with a new block \texttt{rhs} as continuation, and $d$ as a writing location. Then, it recursively compiles both $P$ and $Q$ with a new environment where the name $x$ is bound to $(s, d, 0)$. $P$ is placed on the original block, and $Q$ placed on the new block \texttt{rhs}.

\begin{compilation}{\theta}{\Delta}{\cutc{P}{x: A}{Q}}
{\getendpoints{\theta}{P} + \getendpoints{\theta}{Q}}
    \initsession{s}{\texttt{rhs}}{d}
    \recurse{\theta}{\Delta[x \mapsto (s, d, 0)]}{P}
    \labelir{rhs}
    \recurse{\theta}{\Delta[x \mapsto (s, d, 0)]}{Q}
\end{compilation}

The compilation of a $\parp$ process is shown below. The process starts by generating code to push a new task for executing $Q$ into the task stack. Then, it compiles $P$, which will be executed immediately. AS with the $\cutp$ process, the number of end points of the generated code is the sum of the number of end points of both $P$ and $Q$. Since the task stack is a stack, task pop instructions always execute the most recently pushed task. Thus, it is guaranteed that $Q$ will only execute after $P$ has finished executing completely.

\begin{compilation}{\theta}{\Delta}{\parc{P}{Q}}
{\getendpoints{\theta}{P} + \getendpoints{\theta}{Q}}
    \pushtask{rhs}
    \recurse{\theta}{\Delta}{P}
    \labelir{rhs}
    \recurse{\theta}{\Delta}{Q}
\end{compilation}

The $\emptyp$ process is compiled as shown below. The process marks the end of the current task, as it has no continuation. The flow of control is thus simply passed to the next task in the linked list. Additionally, we mark this instruction as an end point. 

\begin{compilation}{\theta}{\Delta}{\emptyc}{1}
    \poptask{\hasendpoint}
\end{compilation}

\subsection{Close and Wait}

The compilation of a $\closep$ process can be seen below. The session identifier is fetched from the channel environment $\Delta$. This action doesn't have a continuation, as the session is closed and no more operations can be performed on it. The flow of control is simply passed to the continuation of the session, and won't return. Additionally, since there's no continuation, the instruction is marked as an end point.

\begin{compilation}{\theta}{\Delta}{\closec{x}}{1}
    \finishsession{\sessionof{\Delta}{x}}{\hasendpoint}
\end{compilation}

The $\waitp$ process is skipped during compilation: $\getendpoints{\theta}{\waitc{x}{P}} \equiv \getendpoints{\theta}{P}$, $\getinstructions{\theta}{\Delta}{\waitc{x}{P}} \equiv \getinstructions{\theta}{\Delta}{P}$. There's no actual code generated for the wait action, as the action of waiting for the other side of the channel is already performed implicitly by the \ac{SAM} execution scheme, closely mirrored by the continue mechanism we defined earlier.

\subsection{Linear let, scan and print}

The compilation of a $\letp$ process is shown below. The compiler generates the \ac{IR} expression corresponding to the \ac{CLASS} expression $E$, and writes the result to the remote data of the channel $x$. Then, since this process has no continuation, the session is finished, similarly to the $\closep$ process.

\begin{compilation}{\theta}{\Delta}{\letc{x}{E}}
{1}
    \writeexpr{\writelocof{\Delta}{x}}{\getexpression{\Delta}{E}}
    \finishsession{\sessionof{\Delta}{x}}{\hasendpoint}
\end{compilation}

The compilation of a $\scanp$ process is near identical except that instead of computing an expression, we use the \texttt{writeScan} instruction to read from the standard input.

\begin{compilation}{\theta}{\Delta}{\scanc{x}}
{1}
    \writescan{\writelocof{\Delta}{x}}{\slots{\theta}{\tau}} \comment \text{where } x:\tau
    \finishsession{\sessionof{\Delta}{x}}{\hasendpoint}
\end{compilation}

To compile the $\printp$ and $\printlnp$ processes, a single instruction is generated to print the result of the expression $E$ to the standard output. Then, the process continues by recursing on its continuation $P$.

\begin{figure}[H]
    \begin{minipage}{0.45\textwidth}
    \begin{compilationembed}{\theta}{\Delta}{\printc{E}{P}}
    {\getendpoints{\theta}{P}}
        \print{\getexpression{\Delta}{E}}
        \recurse{\theta}{\Delta}{P}
    \end{compilationembed}
    \end{minipage}
    \hfill
    \begin{minipage}{0.45\textwidth}
    \begin{compilationembed}{\theta}{\Delta}{\printlnc{E}{P}}
    {\getendpoints{\theta}{P}}
        \printline{\getexpression{\Delta}{E}}
        \recurse{\theta}{\Delta}{P}
    \end{compilationembed}
    \end{minipage}
\end{figure}

\subsection{Choice and branching}

The compilation of choice processes is shown below. Code is generated which writes a tag to the remote data of the channel $x$. We define $i(\#a)$ as the index of the tag \texttt{a} in the alphabetically sorted list of tags in the offer process. Then, the pseudo-instruction \texttt{continueIfNegative} is inserted, which adds a \texttt{continueSession} instruction if the remainder type of $x$ is negative. Finally, the process recurses on its continuation $P$, advancing the offset of channel $x$ to skip over the written tag, taking into account the remainder type of the channel, as explained in \Cref{subsubsection:AdvancingOffsets}.

\begin{compilation}{\theta}{\Delta}{\choosec{a}{x}{P}}
{\getendpoints{\theta}{P}}
    \writetag{\writelocof{\Delta}{x}}{i(\#a)}
    \continueifneg{x}{A} \comment \text{where } x:A \text{ in }P
    \recurse{\theta}{\advancech{\Delta}{x}{+}{\theta}{\seqst{\tags}}{A}}{P}
\end{compilation}

The compilation of a $\casep$ process is shown below. The compiler generates code which reads a tag from the local data of channel $x$, jumps to the corresponding label and then recurses, generating code for each possible branch, preceded by the respective label. The labels are generated by their alphabetical order, so that their index matches the one written by the \texttt{writeTag} instruction. As before, the offset of channel $x$ is advanced to skip over the read tag. Here, the remainder type of $x$ depends on the chosen branch.

\begin{compilation}{\theta}{\Delta}{\casec{x}{\labelc{a}{A},\,\dots,\,\labelc{z}{Z}}}
    {\max(\getendpoints{\theta}{A},\,\dots,\,\getendpoints{\theta}{Z})}
    \branchtag{\readlocof{\Delta}{x}}{\texttt{a}:\getendpoints{\theta}{A},\,\dots,\,\texttt{z}:\getendpoints{\theta}{Z}}
    \labelir{a}
    \recurse{\theta}{\advancech{\Delta}{x}{-}{\theta}{\seqst{\tags}}{\tau_A}}{A} \comment \text{where } x:\tau_A \text{ in }A
    \beginir\dots
    \labelir{z}
    \recurse{\theta}{\advancech{\Delta}{x}{-}{\theta}{\seqst{\tags}}{\tau_Z}}{Z} \comment \text{where } x:\tau_Z \text{ in }Z
\end{compilation}


\subsection{Send and Receive}

The compilation of a $\sendp$ process is shown below. The process initializes a session $s$ with a new block \texttt{lhs} as continuation and the local data section $d$ as a writing location. $s$ and $d$ are added to the process header, similarly to the $\cutp$ process. In this case, however, $d$ has a layout given by $\alllocalslots{\theta}{A}$, i.e., it only reserves space for received data. This is because, unlike the $\cutp$ process, only this end point of the session will be using this data section to read data.

The initialized session $s$ is then written the remote data of $x$, after which a \texttt{continueIfNegative} instruction is placed. This is necessary, as before, to implement the write-to-read adjustment of the \ac{SAM}. After this, the process recurses immediately on $Q$, advancing the offset of channel $x$ to skip over the sent data, as explained in \Cref{subsubsection:AdvancingOffsets}.

Finally, on the new block \texttt{lhs}, the process recurses on $P$, with a new environment where $y$ is bound to the channel state $(s, d, 0)$. This block will only be executed when the reader of the session we just wrote passes control to it.

\begin{compilation}{\theta}{\Delta}{\sendc{x}{y}{P}{Q}}
{\getendpoints{\theta}{P} + \getendpoints{\theta}{Q}}
    \initsession{s}{\texttt{lhs}}{d}
    \writesession{\writelocof{\Delta}{x}}{s}
    \continueifneg{\sessionof{\Delta}{x}}{A} \comment \text{where } x:A \text{ in }Q
    \recurse{\theta}{\advancech{\Delta}{x}{+}{\theta}{\seqst{\sessions}}{A}}{Q}
    \labelir{lhs}
    \recurse{\theta}{\Delta[y \mapsto (s, d, 0)]}{P}
\end{compilation}

The $\recvp$ process is compiled as shown below. When executing this process, a session has already been written to the local data section associated to channel $x$. A session binding instruction is generated to bind this stored session into a fresh session identifier $s$, which is added to the process header, and a to a fresh local data section $d$, which is also added to the header, with layout $\alllocalslots{\theta}{A}$, similarly to the $\sendp$ process.

The newly bound session must be continued if its type is negative, ensuring that the writing side of the session executes first. After this, the compilation recurses on $P$, with a new environment where $x$ has been advanced to skip over the received session, and $y$ is bound to the channel state $(s, d, 0)$.

\begin{compilation}{\theta}{\Delta}{\recvc{x}{y: A}{P}}
{\getendpoints{\theta}{P}}
    \bindsession{s}{\readlocof{\Delta}{x}}{d}
    \continueifneg{s}{A} \comment \text{where } x:A \text{ in }P
    \recurse{\theta}{\advancech{\Delta}{x}{-}{\theta}{\seqst{\sessions}}{A}[y \mapsto (s, d, 0)]}{P}
\end{compilation}

\subsection{Forwarding}

Given a process $\fwdc{x}{y}$, it can be assumed without loss of generality that $x$ is of negative type and $y$ of positive type - if this is not the case, $x$ and $y$ can be swapped.

Forwarding requires us to move any data which has been received (but not consumed) through $x$ on its local data section to the remote data of $y$, which is positive. This data is known be composed of the slots given by $\slots{\theta}{A}$, where $A$ is the type being forwarded. For example, if $A = \lintt$, then $\seqst{\ints}$ must be moved from the local data of $x$ to the remote data of $y$. This is the first operation performed by the forwarding process, as shown below.

\begin{compilation}{\theta}{\Delta}{\fwdc{x}{y}}
{1}
    \moveslots{\writelocof{\Delta}{y}}{\readlocof{\Delta}{x}}{\slots{\theta}{A}}
    \forwardsessions{\sessionof{\Delta}{x}}{\sessionof{\Delta}{y}}{\hasendpoint}
\end{compilation}

\subsection{Process Calls}

Process calls are compiled as shown below. The identifier $id$ of the \ac{IR} process to call is generated from the name of the process being called and the polarities of its type parameters, using the same scheme as when generating process definitions. This ensures that the correct variant of the process is being called, as explained in \Cref{subsection:ProcessDefinitions}.

\begin{compilation}{\theta}{\Delta}{\proccallgc{processId}{A_1,\,\dots,\,A_n}{x_1,\,\dots,\,x_m;\,y_1,\,\dots,y_k}}{1}
    \callprocess{id}{\hasendpoint}{T, S, D}
\end{compilation}

The list of type arguments $T_1,\,\dots,\,T_n$ is generated from the type parameters $A_1,\,\dots,\,A_n$. Each $T_i$ is the active slot tree of type, computed as $\slots{\theta}{A_i}$.

The list of session arguments $S_1,\,\dots,\,S_m$ is generated from the linear arguments $x_1,\,\dots,\,x_m$. Each $S_i$ is a tuple $(s_t, s_s, o)$, where $s_t$ is the session identifier of the argument in the environment of the process being called, $s_s = \sessionof{\Delta}{x}$, and $o = \offsetof{\Delta}{x}$.

The list of data arguments $D_1,\,\dots,\,D_{m+k}$ is generated similarly to the session arguments, but with extra entries for the exponential arguments. Each $D_i$ is a tuple $(d_t, d_s, t, c)$.

For $i \leq m$ (the linear arguments), $d_t$ is the target data section identifier of the argument in the environment of the process being called, $d_s = \readlocof{\Delta}{x}$, $t = \slots{\theta}{\tau}$ where $x_i: \tau$ and $c$ is the flag indicating whether the argument data should be cloned or moved, which in this case is always \texttt{move}, as $x_i$ is linear.

For $i > m$ (the exponential arguments), let $j = i - m + 1$. $d_t$ is once again the target data section identifier of the argument in the environment of the process being called, $d_s = \readlocof{\Delta}{y_j\textsc{}}$, $t = \slots{\theta}{\whynott{\tau}}$ where $y_j: \tau$. In this case, $c$ is always \texttt{clone}, as $y_j$ is exponential.

\todo{this is quite complex, maybe this could be written in a clearer way?}

\todo{an example of a process call could be nice here}

\subsection{Creating and calling exponentials}

The right-hand-side $P$ of a $\bangc{x}{y}{P}$ process is compiled as a separate \ac{IR} process, with a unique identifier generated from the identifier of the process definition being compiled. This new process inherits the same type variables as the current process. The argument session $y$ is bound to a new session identifier $s$ in the new process header, along with a new local data section $d$ with layout $\alllocalslots{\theta}{A}$, where $A$ is the type of channel $y$. The number of end points of the new process' header is set to $\getendpoints{\theta}{P}$.

All free names in $P$ which are not $y$ are exponentials (guaranteed by the typing rules of \ac{CLASS}), and should be passed to the new process as arguments. For each of them, a data section is created in the new process header, with layout $\alllocalslots{\theta}{\whynott{\tau}}$, where $\tau$ is the type of the free name. The environment used to compile $P$ is also updated to bind these names to their respective data sections, with offset $0$.

On the original process, an exponential write instruction is generated referencing the new process identifier, along with all necessary type arguments (derived from the type parameters of the current process), and all data arguments for each of the free names in $P$. Finally, since the $\bangp$ process has no continuation, the session is finished, similarly to the $\closep$ process.

\begin{compilation}{\theta}{\Delta}{\bangc{x}{y}{P}}
{1}
    \writeexp{\writelocof{\Delta}{x}}{id}{T, D}
    \finishsession{\sessionof{\Delta}{x}}{\hasendpoint}
\end{compilation}

The list of type arguments $T_1,\,\dots,\,T_n$ is generated similarly to the type arguments of a process call, but inheriting the same exact type parameters as the current process.

The list of data arguments $D_1,\,\dots,\,D_m$ is generated from the free names in $P$, excluding $y$. Each $D_i$ is a tuple $(d_t, d_s, t, c)$, where $d_t$ is the target data section identifier of the argument in the environment of the process being called, $d_s = \readlocof{\Delta}{x_i}$, $t = \slots{\theta}{\whynott{\tau}}$ where $x_i: \tau$ and $c$ is the flag indicating whether the argument data should be cloned or moved, which in this case is always \texttt{clone}, as all free names in $P$ are exponentials.

The compilation of a $\whynotc{x}P$ process is shown below. A new drop bit $(S, d)$ identified by a fresh drop bit identifier $e$ is inserted into the process header, with slots $S = \slots{\theta}{A}$, where $A$ is the type of channel $x$, and $d = \readlocof{\Delta}{x}$, the local data location of channel $x$, where the exponential data with slots $S$ is stored.
The compilation of a $\whynotc{x}P$ process is shown below. A new drop bit $(S, d)$ identified by a fresh drop bit identifier $e$ is inserted into the process header, with slots $S = \slots{\theta}{A}$, where $A$ is the type of channel $x$, and $d = \readlocof{\Delta}{x}$, the local data location of channel $x$, where the exponential data with slots $S$ is stored.

Then, a defer drop instruction is generated, which sets the drop bit $e$ to true, ensuring that the exponential data associated to it is dropped when the process ends. Finally, the compilation recurses on $P$, with no changes to the environment, as no changes to the channel state are made.

\begin{compilation}{\theta}{\Delta}{\whynotc{x}{P}}{\getendpoints{\theta}{P}}
    \deferdrop{e}
    \recurse{\theta}{\Delta}{P}
\end{compilation}

\subsection{Unfolds}

Unfolded processes are compiled such that positive unfolds ($\unfoldmup$) always pass control to their continuation. If the remainder type is negative, i.e., the next action will read data from the channel, then we must pass control to the other side so that it can write the data before we read it. If the remainder type is positive, we cannot immediately continue as we must first give an opportunity for the other side to read any data we may have already written relative to the previous iteration of the unfold.

Negative unfolds ($\unfoldnup$) only pass control to their continuation if the remainder type is negative, as in this case, the positive side must have already executed a positive unfold, and thus, given control to the negative side. If the remainder type is negative, the positive side must execute first, and thus we surrender control back to it. If the remainder type is positive, we can continue executing immediately.

In both cases, the process finishes by recursing on the right-hand-side $P$ of the unfold, with a new environment where $x$ is bound to a new channel state equal to its previous state, but with offset $0$, as a new segment of the session type starts.

\begin{figure}[H]
    \centering
    \begin{minipage}{0.45\textwidth}
        \begin{compilationembed}{\theta}{\Delta}{\unfoldmuc{x}{P}}
        {\getendpoints{\theta}{P}}
            \continuesession{\sessionof{\Delta}{x},\,\texttt{cont}}
            \labelir{cont} 
            \recurse{\theta}{\Delta[x \mapsto (\sessionof{\Delta}{x}, \localdataof{\Delta}{x}, 0)]}{P}
        \end{compilationembed}
    \end{minipage}
    \hfill
    \begin{minipage}{0.45\textwidth}
        \begin{compilationembed}{\theta}{\Delta}{\unfoldnuc{x}{P}}
        {\getendpoints{\theta}{P}}
            \beginir\%\ \text{where } x:A \text{ in }P
            \continueifneg{x}{A}
            \recurse{\theta}{\Delta[x \mapsto (\sessionof{\Delta}{x}, \localdataof{\Delta}{x}, 0)]}{P}
        \end{compilationembed}
    \end{minipage}
\end{figure}


\subsection{Receiving and sending types}

The \ac{IR} only allows polymorphic data using type variables which are part of the process header. Thus, to implement the $\recvtyc{x}{X}P$ process, a new process must be called which has $X$ as a type parameter, and depending on $X$'s polarity. This mechanism is similar to the one used for exponentials, but in this case, the process is immediately called instead of writing an exponential, and it may have free linear names which must be passed as arguments. Additionally, two processes must be generated, one for each polarity of $X$.

The compilation definition of the $\recvtyc{x}{X}P$ process is shown below. Two new process identifiers $id^-$ and $id^+$ are generated from the identifier of the current process. For each of them, a new process is generated with the same type parameters as the current process, plus a new type parameter $X$ with polarity negative for $id^-$ and positive for $id^+$. The number of end points of both new processes is set to $\getendpoints{\theta}{P}$.

First, a branch instruction is generated which reads a tag from the local data of channel $x$, which indicates the polarity of the type being received. Then, depending on the read tag, either $id^-$ or $id^+$ is called, passing all free names in $P$ as arguments, along with the new session which represents $x$ in $P$, stored in $\readlocof{\Delta}{x} + \slotoffset{\seqst{\types}}{\seqst{\sessions; \tags}}$. Both $T^-$ and $T^+$ are generated from the type parameters of the current process, exactly as in a process call, but with an extra entry for $X$, which is read from $\readlocof{\Delta}{x}$.

All session arguments $S$ and data arguments $D$ are generated from the free names in $P$, simply forwarding their channel states from the environment $\Delta$. \todo{do I need to clarify this?}

\begin{compilation}{\theta}{\Delta}{\recvtyc{x}{X}P}{1}
    \branchtag{\concatst{\readlocof{\Delta}{x}}{\slotoffset{\seqst{\types; \sessions}}}{\seqst{\tags}}}{\texttt{neg:1},\,\texttt{pos:1}}
    \labelir{neg}
    \callprocess{id^-}{\hasendpoint}{T^-, S, D}
    \labelir{pos}
    \callprocess{id^+}{\hasendpoint}{T^+, S, D}
\end{compilation}

On the other hand, the $\sendtyc{x}{A}{P}$ process is compiled in a more direct manner, as shown below. This is possible because the type variable has a concrete type $A$ at the point of sending, known at compile time.

First, a fresh session $s$ is initialized with a new block \texttt{cont} as continuation, and a fresh local data section $d$ with layout $\alllocalslots{\theta}{B}$, where $B$ is the type of channel $x$ in $P$. Then, the type $A$ is written to the remote data of channel $x$, along with a $\sessions$ referencing a $s$ and a $\tags$ indicating $A$'s polarity ($0$ for negative types, $1$ for positive types). The session of channel $x$ is then finished.

The right-hand-side $P$ is compiled on the new block \texttt{cont}, with a new environment where $x$ is rebound to the channel state $(s, d, 0)$. This block will be executed by the $\recvtyp$ process, similarly to how the right-hand-side of a $\sendp$ process is executed by $\recvp$.

\begin{compilation}{\theta}{\Delta}{\sendtyc{x}{A}{P}}
{\getendpoints{\theta}{P}}
    \initsession{s}{\texttt{cont}}{d}
    \writetype{\writelocof{\Delta}{x}}{\slots{\theta}{A}}
    \writesession{\writelocof{\Delta}{x} + \slotoffset{\seqst{\types}}{\seqst{\sessions; \tags}}}{s}
    \writetag{\writelocof{\Delta}{x} + \slotoffset{\seqst{\types; \sessions}}{\seqst{\tags}}}{p} \comment \text{where } p = 0 \text{ if } A \text{ is negative, } 1 \text{ otherwise}
    \finishsession{\sessionof{\Delta}{x}}{\hasendpoint}
    \labelir{cont}
    \recurse{\theta}{\Delta[x \mapsto (s, d, 0)]}{P}
\end{compilation}

\subsection{Affine, use and discard}

The compilation of the $\affinec{x}P$ process, shown below, implements the affine protocol described in \Cref{subsection:IRDroppingData}: a tag is received from $x$ (with an offset of 0) indicating whether the channel is being used or discarded. If it's being used, the process continues with an environment where $x$ has been advanced by the read tag. If the tag indicates the session is being discarded, any co-affine data free in $P$ is discarded and the affine session is finished.

One important aspect to note is that although $\affinet{A}$ is a positive type, its process starts by reading a tag, while the dual type $\coaffinet{\dualt{A}}$, which is negative, starts by writing a tag. The rest of the compilation scheme guarantees that the positive side always executes first, which in this case is the reading side. Thus, the $\affinep$ process must always pass control to the negative side before continuing.

An alternative solution would be to change the definition of $\affinet{A}$ to be negative and $\coaffinet{\dualt{A}}$ to be positive, but this would add a lot of complexity the important affine value optimization described in \Cref{subsection:AffineValueOptimizations}.

\begin{compilation}{\theta}{\Delta}{\affinec{x}{P}}
{\getendpoints{\theta}{P}}
    \continuesession{\sessionof{\Delta}{x},\,\texttt{cont}}
    \labelir{cont}
    \branchtag{\localdataof{\Delta}{x}}{\texttt{use:}\getendpoints{\theta}{P},\,\texttt{discard:1}}
    \labelir{use}
    \recurse{\theta}{\advancech{\Delta}{x}{-}{\theta}{\seqst{\tags}}{A}}{P} \comment \text{where } x:A \text{ in }P
    \labelir{discard}
    \recurse{\theta}{\Delta}{\discardc{y_i}} \comment \text{for each free co-affine name }y_i\text{ in }P
    \finishsession{\sessionof{\Delta}{x}}{\hasendpoint}
\end{compilation}

To compile a $\discardc{x}$ process, as seen below, a tag is written to the remote data of channel $x$ indicating that the session is being discarded. Then, control is passed to the continuation of the session. When it returns, the current task is finished, unless the $\discardp$ process was generate as part of the discard branch of an affine process, in which case the execution of the current task continues.

\begin{compilation}{\theta}{\Delta}{\discardc{x}}
{1}
    \writetag{\writelocof{\Delta}{x}}{1} \comment \text{indicating discard}
    \continuesession{\sessionof{\Delta}{x},\,\texttt{cont}}
    \labelir{cont}
    \poptask{\hasendpoint} \comment \text{unless part of an affine discard branch}
\end{compilation}

Finally, the $\usec{x}{P}$ process is compiled as shown below. A tag is written to the remote data of channel $x$ indicating that the session is being used. Then, if the continuation is negative, the control is passed to the other side of the session. Otherwise, the process continues executing immediately.

\begin{compilation}{\theta}{\Delta}{\usec{x}{P}}
{\getendpoints{\theta}{P}}
    \writetag{\writelocof{\Delta}{x}}{0} \comment \text{indicating use}
    \continueifneg{x}{A} \comment \text{where } x:A \text{ in }P
    \recurse{\theta}{\advancech{\Delta}{x}{+}{\theta}{\seqst{\tags}}{A}}{P}
\end{compilation}

\subsection{Shared state}

The compilation of the $\cellc{x}{y}{P}$ process, which creates a new shared cell on $x$, is shown below. First, a new cell is written to the remote data of channel $x$, with space for a single $\sessions$ slot. Then, a new session $s$ is initialized with a new block \texttt{cont} as continuation, and a fresh local data section $d$ with layout $\alllocalslots{\theta}{A}$, where $A$ is the type of channel $y$ in $P$. This session $s$ is written to the cell we just created, using the cell as a base address, and an offset of $0$. Finally, the process recurses on $P$, with a new environment where $y$ is bound to the channel state $(s, d, 0)$. This block will be executed by the process which takes the cell's content.

\begin{compilation}{\theta}{\Delta}{\cellc{x}{y}{P}}
{\getendpoints{\theta}{P}}
    \writecell{\writelocof{\Delta}{x}}{\seqst{\sessions}}
    \initsession{s}{\texttt{cont}}{d}
    \writesession{c(\writelocof{\Delta}{x})}{s}
    \labelir{cont}
    \recurse{\theta}{\Delta[y \mapsto (s, d, 0)]}{P}
\end{compilation}

The $\takec{x}{y}P$ process is compiled as shown below. First, the cell's mutex is locked. Then, the session stored in the cell is bound to a fresh session identifier $s$, which is added to the process header, along with a fresh local data section $d$ with layout $\alllocalslots{\theta}{A}$, where $A$ is the type of channel $y$ in $P$. If the type $A$ is negative, the session is continued, similarly to the $\recvp$ process. Finally, the process recurses on $P$, with a new environment where $y$ is bound to the channel state $(s, d, 0)$. The cell's mutex is unlocked when the session $s$ is released.

\begin{compilation}{\theta}{\Delta}{\takec{x}{y}{P}}
{\getendpoints{\theta}{P}}
    \lockcell{\readlocof{\Delta}{x}}
    \bindsession{s}{c(\readlocof{\Delta}{x})}{d}
    \continueifneg{s}{A} \comment \text{where } x:A \text{ in }P
    \recurse{\theta}{\Delta[y \mapsto (s, d, 0)]}{P}
\end{compilation}

The $\putc{x}{y}{P}Q$ process is compiled as shown below. First, a fresh session $s$ is initialized with a new block \texttt{lhs} as continuation, and a fresh local data section $d$ with layout $\alllocalslots{\theta}{A}$, where $A$ is the type of channel $y$ in $P$. This session $s$ is written to the cell referenced by channel $x$. Then, the cell is unlocked and the process recurses immediately on $Q$. Finally, on the new block \texttt{lhs}, the process recurses on $P$, with a new environment where $y$ is bound to the channel state $(s, d, 0)$. This block will only be executed when the taker of the cell we just wrote passes control to it.

\begin{compilation}{\theta}{\Delta}{\putc{x}{y}{P}{Q}}
{\getendpoints{\theta}{P} + \getendpoints{\theta}{Q}}
    \initsession{s}{\texttt{lhs}}{d}
    \writesession{c(\readlocof{\Delta}{x})}{s}
    \unlockcell{\readlocof{\Delta}{x}}
    \recurse{\theta}{\Delta}{Q}
    \labelir{lhs}
    \recurse{\theta}{\Delta[y \mapsto (s, d, 0)]}{P}
\end{compilation}

The process $\releasec{x}$ is compiled directly to the instruction which decrements the reference count of the cell referenced by channel $x$, as shown below. The responsibility of correctly dropping the cell's contents when the reference count is zero is left to the \ac{IR} implementation (the C code generation phase), as explained in \Cref{subsubsection:CellSlotArguments}.

\begin{compilation}{\theta}{\Delta}{\releasec{x}}{1}
    \releasecell{\readlocof{\Delta}{x}}
    \poptask{\hasendpoint}
\end{compilation}

Finally, the $\sharec{x}{P}{Q}$ process is compiled as shown below. Since the cell is going to be shared between two processes, its reference count must be incremented. By default, with concurrency enabled, a new thread is launched to execute the right-hand-side $Q$ of the share process, as shown below. If it is disabled, the right-hand-side is instead added as a new task to the task queue, which will be executed when the current task finishes.

\begin{compilation}{\theta}{\Delta}{\sharec{x}{P}{Q}}
{\getendpoints{\theta}{P} + \getendpoints{\theta}{Q}}
    \acquirecell{\readlocof{\Delta}{x}}
    \launchthread{\texttt{rhs}}
    \recurse{\theta}{\Delta}{P}
    \labelir{rhs}
    \recurse{\theta}{\Delta}{Q}
\end{compilation}

\section{Generating C code from IR programs}

The last phase of the compilation pipeline is the C generation phase. This phase takes as input the \ac{IR} produced by the compilation scheme described in the previous section, and produces a C program which implements the semantics of the \ac{IR}.

This C program can then be compiled with \ac{GCC} to produce an executable which runs a computation equivalent to the original \ac{CLASS} program. The generated code depends only on the standard C library and on \texttt{pthread}, if concurrency is enabled. Additionally, it makes extensive use of the \ac{GCC} extension for computed gotos, which allows for efficient implementation of continuations.

The first subsection of this section describes how the slot trees are translated to concrete data structures in C. The second subsection describes the environment data structure, which stores the state of a process invocation.  The third subsection shows how the main function and the generated code are laid out. Finally, the the last subsections give an overview of how the main instructions of the \ac{IR} abstract machine are implemented in C.

\subsection{Slot tree memory layouts}

In the \ac{IR} abstract machine, the concept of data is abstracted away into slots, covered in \Cref{subsection:IRSlots}. A slot tree describes the possible slots which may be stored in a data section in sequence, with branching. A set of slot trees describes the multiple possible slot trees which may be stored in a data section, with only one of them being active at any given time.

To generate C code, these slot trees must be translated to concrete, byte-addressed, memory layouts. To this end, a memory layout is defined as a tuple $L = (S, A)$, where $S$ is the size in bytes of the layout, and $A$ is the alignment in bytes the layout must have in memory. Both $S$ and $A$ are not constants, but C expressions which can be evaluated at runtime. This is necessary because polymorphic types have layouts which are only known at runtime.

In order to satisfy the alignment requirements of the various C types used to implement the \ac{IR} abstract machine, a C macro \texttt{ALIGN(s, a)} is defined, which rounds up the size \texttt{s} to the nearest multiple of the alignment \texttt{a}. This macro is used extensively when computing memory layouts. Additionally, a C macro \texttt{MAX(a, b)} is defined, which computes the maximum of two integers \texttt{a} and \texttt{b}.

Two memory layouts $L_1 = (S_1, A_1), L_2 = (S_2, A_2)$ can be concatenated to form a new layout $L = L_1 + L_2$, where $L = (\texttt{ALIGN}(S_1, A_2) + S_2,\ \texttt{MAX}(A_1, A_2))$, such that $L_1$ is placed first in memory, followed by $L_2$, with padding added to ensure that $L_2$ starts at an address which is a multiple of $A_2$. Additionally, the maximum of two layouts $L = L_1 \vee L_2$ is defined as $L = (\texttt{MAX}(S_1, S_2),\ \texttt{MAX}(A_1, A_2))$.

The memory layout of a slot $s$ on a given target architecture $\arch$ is given by the function $\slotlayout{\sigma}{s}$ (defined in \Cref{fig:CSlotLayout}), where $\sigma$ is a map between type variables and their layouts, and $\arch$ a target architecture. A target architecture $\arch$ is a mapping from C types to their memory layouts, such that $\archlayout{type} = (S, A)$ where $S$ and $A$ where, respectively, the size and alignment of \texttt{type} on the target architecture. A concrete definition of $\sigma$ will be introduced in \Cref{subsection:CEnvironments}.

\begin{figure}[H]
    \begin{align*}
        &\begin{aligned}
            &\slotlayout{\sigma}{\ints} = \archlayout{int} && \slotlayout{\sigma}{\bools} = \archlayout{unsigned char} \\
            &\slotlayout{\sigma}{\strings} = \archlayout{char*} && \slotlayout{\sigma}{\sessions} = \archlayout{struct session*} \\
            &\slotlayout{\sigma}{\tags} = \archlayout{unsigned char} && \slotlayout{\sigma}{\types} = \archlayout{struct type} \\
            &\slotlayout{\sigma}{\exponentials} = \archlayout{struct exponential*} && \slotlayout{\sigma}{\cells} = \archlayout{struct cell*} \\
        \end{aligned}\\
        &\slotlayout{\sigma}{t_i} = \layoutofvar{\sigma}{t_i}\\
    \end{align*}
    \caption{Memory layout of a slot $s$ on a target architecture $\arch$ and a type variable layout map $\sigma$.}
    \label{fig:CSlotLayout}
\end{figure}

The definition of $\slotlayout{\sigma}{s}$ is extended to support both trees of slots, and sets of slot trees, as shown in \Cref{fig:CTreeLayout}. The layout of a tree of slots is defined as the concatenation of the layouts of each of its slots, and maximum across any branches, while the layout of a set of slot trees is defined as the maximum layout among all its trees.

\begin{figure}[H]
    \begin{align*}
        &\slotlayout{\sigma}{\nilst} = (0, 1)\\
        &\slotlayout{\sigma}{\seqst{s; T}} = \slotlayout{\sigma}{s} + \slotlayout{\sigma}{T}\\
        &\slotlayout{\sigma}{\tagst{T_1 \tagsepst \dots \tagsepst T_n}} = \bigvee_{i=1}^{n} \slotlayout{\sigma}{T_i}\\
        &\slotlayout{\sigma}{\left\{T_1,\,\dots,\,T_n\right\}} = \bigvee_{i=1}^{n} \slotlayout{\sigma}{T_i}
    \end{align*}
    \caption{Memory layout of slot trees $T$ and of a set of slot trees on a target architecture $\arch$ and a type variable layout map $\sigma$.}
    \label{fig:CTreeLayout}
\end{figure}

\todo{might need to also define 'current layout' to get the current size of a branching tree}

\subsection{Environments} \label{subsection:CEnvironments}

The main data structure used to implement the execution state of the \ac{IR} abstract machine is the environment, which stores all state related to a process invocation. In the generated code, environments are compiled to heap allocated byte arrays, which are indexed through an environment layout $\envlayout$.

An environment layout $\envlayout$ is a compile-time data structure which describes the byte offsets at which each entry of the environment can be found, and is computed from the header of the \ac{IR} process being compiled. The address of the end point counter (a C integer or atomic integer if concurrency is enabled) is given by $\envlayouteps{\texttt{env}}$ where \texttt{env} is a pointer to the environment. Drop bits are stored at $\envlayoutdrop{\texttt{env}}$, type information for type variable $t_i$ is stored at $\envlayouttype{\texttt{env}}{t_i}$, sessions $s_i$ are stored at $\envlayoutsession{\texttt{env}}{s_i}$, and local data $d_i$ at $\envlayoutlocal{\texttt{env}}{d_i}$.

The layout of the environment is computed trivially by concatenating the layouts of each of its entries, in the order they appear in the process header. The layout for for data entries is the most noteworthy, as it depends on the stored data type. This layout is computed from $\slotlayout{\sigma}{S}$, where $S$ is the set of slot trees of the data being stored, and $\sigma$ is a mapping between type variables and their layouts, fetched from the environment itself (e.g., $\sigma(t_i) = \texttt{access}(\envlayouttype{\texttt{env}}{t_i})$). Its formal definition is omitted for brevity.

Sessions are stored in the environment as instances of the C type \texttt{struct session}, shown below. Sessions store the address of a label to jump to the continuation (\texttt{cont}), a pointer to the environment to switch to before jumping to the continuation (\texttt{cont\_env}), a pointer to the remote data (writing buffer) of the session (\texttt{cont\_data}), and a pointer to the session entry in the continuation environment pertaining to the other end point of the session (\texttt{cont\_session}).

\begin{lstlisting}[language=C]
struct session {
  void* cont;
  char* cont_env;
  char* cont_data;
  struct session* cont_session;
};
\end{lstlisting}

Type information is stored in the environment as instances of the C type \texttt{struct type}, shown below. This structure holds the \texttt{size} and \texttt{alignment} of the data described by the type. This structure is also used to store the value of $\types$ slots in the local data sections.

\begin{lstlisting}[language=C]
struct type {
    int size;
    int alignment;
};
\end{lstlisting}

\subsubsection{Cleaning up an environment} \label{subsubsection:CEnvironmentCleanup}

When the end point counter of an environment reaches zero, any set drop bits must be processed, dropping their associated data. Thus, whenever an instruction decrements the end point counter, it must check if it reached zero, and if so, run the cleanup code. Only then can the environment itself be freed. For example, the code below decrements the end point counter, checking if it reached zero, and if so, checks if the drop bit $0$ is set, dropping its associated data if so, before freeing the environment itself.

\begin{lstlisting}[language=C]
if (atomic_fetch_sub(&(*(atomic_int*)(env)), 1) == 1) {
    if ((*(unsigned char*)(env + 4)) & (1 << 0)) {
        /* drop data associated to drop bit 0 */
    }
    /* other drop bit checks... */
    managed_free(env);
}
\end{lstlisting}

If a process has only one end point, a small optimization is applied which removes the end point counter from the environment, and any decrement instructions are replaced with the cleanup code directly, as the end point counter is guaranteed to reach zero when these instructions are reached.

\subsection{Execution and registers}

The \texttt{main} function of the generated C program initializes any required global state (such as the allocator, or profiling tools), and invokes a function named \texttt{executor} with a \texttt{NULL} argument. The \texttt{executor} function is the core of the \ac{IR} abstract machine implementation. Each thread of execution runs its own instance of this function, receiving as argument a pointer to its entry task.

At the start of the \texttt{executor} function, registers used to implement the \ac{IR} instructions are declared. The first two registers \texttt{task} and \texttt{env} respectively store pointers to the head of the task stack (introduced in the next subsection) and the current environment, which is state kept across \ac{IR} instructions. The rest of the registers are used to store temporary values used within the execution of each instruction. After declaring the registers, the thread state is initialized.

\begin{lstlisting}[language=C]
void executor(struct task* entry) {
    /* Registers */
    register struct task* task;
    register char* env;
    register void* tmp_ptr1;
    register void* tmp_ptr2;
    register int tmp_int;
    struct session tmp_session;

    /* Initialize thread state (...) */
\end{lstlisting}

After initializing the thread state, control is passed to the entry task (received as argument), which contains a continuation address and environment. If the entry task is \texttt{NULL}, the main process is called instead, as shown below. To call the main process, a new environment is allocated with the size computed from its header, and a jump is performed to its entry block.

\begin{lstlisting}[language=C]
    if (entry == NULL) {
        /* callProcess(main) */
        tmp_ptr1 = managed_alloc(144);
        env = ((char*)(tmp_ptr1));
        goto main_entry;
    } else {
        env = entry->cont_env;
        tmp_ptr1 = entry->cont;
        managed_free(entry);
        goto *tmp_ptr1;
    }
\end{lstlisting}

After these instructions, still within the \texttt{executor} function, for each block of each process a unique C label is generated, followed by the C code generated from the \ac{IR} instructions in the block. Finally, a label named \texttt{end} is generated, which represents the end of the thread's execution.

\begin{lstlisting}[language=C]
main_entry:
    /* instructions of main's entry block ... */
foo_cut_lhs_0:
    /* instructions of foo's cut_lhs_0 block ... */
    /* other blocks ... */
end:
    return 0;
}
\end{lstlisting}

\subsection{Tasks}

Tasks are stored in a linked list local to each thread of execution, where each node is of the type shown below. The \texttt{task} register points to the head node of this stack. The \texttt{next} field points to the next task to be executed, while the \texttt{cont} and \texttt{cont\_env} fields store the program address and environment to switch to when executing the task. 

\begin{lstlisting}[language=C]
struct task {
    struct task* next;
    void* cont;
    char* cont_env;
};
\end{lstlisting}

The \texttt{task} register is initialized at the start of the \texttt{executor} function with a dummy task which jumps to the \texttt{end} label, ending the thread's execution. This dummy task is necessary to ensure that the task stack is never empty, simplifying the implementation of the \texttt{popTask} instruction.

\begin{lstlisting}[language=C]
    task = managed_alloc(sizeof(struct task));
    task->cont = &&end;
\end{lstlisting}

The two \ac{IR} instructions for manipulating tasks, \texttt{pushTask} and \texttt{popTask}, manipulate the linked list of tasks pointed to by the \texttt{task} register. The first instruction, \texttt{pushTask}, allocates a new task, sets its continuation to the C label of the chosen block, and inserts it at the head of the linked list. The generated C code for this instruction is shown below.

\begin{lstlisting}[language=C]
    /* pushTask(cont) */
    tmp_ptr1 = task;
    task = managed_alloc(sizeof(struct task));
    task->next = ((struct task*)(tmp_ptr1));
    task->cont = &&proc_cont; /* proc_cont is the C label of block 'cont' */
    task->cont_env = env;
\end{lstlisting}

The \texttt{popTask} instruction, whose generation is shown below, removes the first task from the linked list, changes the current environment and jumps to its continuation.

\begin{lstlisting}[language=C]
    /* popTask() */
    tmp_ptr1 = task;
    tmp_ptr2 = task->cont;
    env = task->cont_env;
    task = task->next;
    managed_free(tmp_ptr1);
    goto *tmp_ptr2;
\end{lstlisting}

\subsection{Session manipulation}

The \texttt{initializeSession} instruction is compiled as shown below. The session at $\envlayoutsession{\texttt{env}}{s_0}$ (in this example, $\texttt{env} + 8$) is initialized with the given continuation label, the current environment as its continuation environment, the address $\envlayoutlocal{\texttt{env}}{d_0}$ (in this case, $\texttt{env} + 40$) as its continuation data, and a pointer to itself as its continuation session.

\begin{lstlisting}[language=C]
    /* initializeSession(s0, cut_lhs_0, d0) */
    (*(struct session*)(env + 8)) = (struct session){
        .cont=&&main_cut_lhs_0,
        .cont_env=env,
        .cont_data=env + 40,
        .cont_session=env + 8
    };
\end{lstlisting}

The compilation of the \texttt{finishSession} instruction is shown below. First, the remote session of the finished session is unlinked from the current session by setting its \texttt{cont\_session} field to itself. This is necessary to avoid use after free errors in instructions such as \texttt{forwardSessions} and \texttt{callProcess}. Then, the continuation is stored on the registers and the environment's end point counter is decremented atomically (when concurrency is disabled, a normal integer is used instead). If the counter reaches zero, the environment cleanup code described in \Cref{subsubsection:CEnvironmentCleanup} is executed. Finally, the environment and program address are switched to the continuation values, and a jump is performed.

\begin{lstlisting}[language=C]
    /* finishSession(s0, end point) */
    (*(*(struct session*)(env + 8)).cont_session).cont_session = (*(struct session*)(env + 8)).cont_session;
    tmp_ptr1 = (*(struct session*)(env + 8)).cont;
    tmp_ptr2 = (*(struct session*)(env + 8)).cont_env;
    if (atomic_fetch_sub((atomic_int*)(env), 1) == 1) {
        /* environment cleanup code */
    }
    env = ((char*)(tmp_ptr2));
    goto *tmp_ptr1;
\end{lstlisting}

The \texttt{continueSession} instruction is compiled similarly to \texttt{finishSession}, as shown below, except that the end point counter is not decremented, and that the continuation of the remote session must be updated so that it eventually jumps back to the correct label.

\begin{lstlisting}[language=C]
    /* continueSession(s0, continue_12) */
    tmp_ptr1 = (*(struct session*)(env + 8)).cont;
    (*(*(struct session*)(env + 8)).cont_session).cont = &&main_continue_12;
    env = (*(struct session*)(env + 8)).cont_env;
    goto *tmp_ptr1;
\end{lstlisting}

The \texttt{forwardSessions} instruction first stores the continuation of its second argument in the temporary registers. Then, it sets the remote sessions of both arguments to point to each other, ensuring no references are left to the local sessions. A temporary variable is used to store the first argument's session, as its value may be overwritten due to the possibility of a session being its own remote session. Finally, the end point counter is decremented, and control is passed to the continuation stored in the temporary registers.

\begin{lstlisting}[language=C]
    /* forwardSessions(s0, s1, end point) */
    tmp_ptr1 = (*(struct session*)(env + 40)).cont;
    tmp_ptr2 = (*(struct session*)(env + 40)).cont_env;
    tmp_session = *(struct session*)(env + 8);
    *(*(struct session*)(env + 8)).cont_session = *(struct session*)(env + 40);
    *(*(struct session*)(env + 40)).cont_session = tmp_session;
    if (atomic_fetch_sub(*(atomic_int*)(env), 1) == 1) {
        /* environment cleanup code */
    }
    env = (char*)tmp_ptr2;
    goto *tmp_ptr1;
\end{lstlisting}

\subsection{Data manipulation}

\todo{Explain moveSlots}

\todo{Explain cloneSlots}

\todo{Explain dropSlots}

\subsection{Process calls}

\todo{rewrite everything here}

Process calls are performed through the \texttt{call} instruction. This instruction receives the label of the process to be called, the size of the environment to be created, and a mapping containing the indices of the session records to be passed as arguments. The generated C code for this instruction is shown below.

\begin{lstlisting}[language=C]
    /* call(proc, size, args) */
    tmp_env = malloc(sizeof(struct environment) + size * sizeof(struct record*));
    /* for each i in 0..#args */
    tmp_env->records[/*i*/] = env->records[args[/*i*/]];
    env = tmp_env;
    goto proc;
\end{lstlisting}

\subsection{Exponentials}

\todo{Write this section}

\subsection{Cell operations}

\todo{Write this section}
